%\section*{Declaración del Uso de IA Generativa}
%
% Actualizada el 31-ago-2026 por indicacion del autor.

\noindent\textbf{Herramienta/s utilizada/s:}

\vspace{0.5em}
Claude (Anthropic), utilizado a través de Claude Code sobre el repositorio del
proyecto, desde las fases iniciales hasta la redacción de esta memoria (marzo a
agosto de 2026).

\vspace{2em}
\noindent\textbf{Uso realizado (marcar lo que corresponda):}
% Para realizar la selección de una casilla inferior, sustituya $\square$ por $\boxtimes$

$\boxtimes$ Generación de texto

$\boxtimes$ Reformulación

$\square$ Traducción / corrección

$\boxtimes$ Sugerencia de estructura

$\square$ Apoyo metodológico

$\boxtimes$ Ayuda para codificar / programar

$\boxtimes$ Otros: generación del material gráfico

\vspace{1em}
\noindent\textbf{Descripción del uso realizado:}

\vspace{0.5em}
\sloppy
La herramienta se utilizó a lo largo de todo el proyecto en tres tareas. La
primera es la escritura y depuración de los programas: los de extracción y
vectorización del tráfico, los de evaluación y las implementaciones de las
arquitecturas de aprendizaje profundo, tanto las propias (\texttt{ByteCNN},
\texttt{ByteLSTM} y el híbrido de dos ramas) como las tomadas del estado del
arte para contrastarlas (\texttt{ByteTransformer} y
\texttt{CrossAttentionNet}). La segunda es la generación del material gráfico:
las figuras y las tablas de esta memoria se producen con programas escritos con
esa asistencia. La tercera es la redacción, estructuración y reformulación de
partes de la memoria y de la documentación técnica del repositorio.

\vspace{0.5em}
\noindent\textit{Dirección y decisiones.} La dirección del trabajo fue en todo
momento del autor: el planteamiento del problema y la tesis central, la
elección de los días del conjunto de datos a procesar, el orden de prioridades,
el alcance de cada experimento y las decisiones sobre qué líneas abandonar o
profundizar. Varias correcciones de rumbo, incluida la detección de resultados
mal interpretados y la petición de contrastar hipótesis que resultaron falsas,
partieron del autor.

\vspace{0.5em}
\noindent\textit{Trazabilidad, y sus límites.} Las contribuciones asistidas
están identificadas en el historial de control de versiones mediante la marca
\texttt{Co-Authored-By}. Esa marca se empezó a aplicar de forma sistemática con
el proyecto ya avanzado, de modo que constituye una cota inferior y no un
registro completo: hubo asistencia también en aportaciones anteriores que no la
llevan. Esta declaración prevalece sobre el recuento del historial.

\vspace{2em}
\noindent\textbf{Compromiso de autoría:}\\
Confirmo que todo el contenido generado ha sido revisado, reelaborado y validado
personalmente.
